\PBchapter{Parede Branca}
\cleardoublepage
Um dos grandes valores para mim é a honestidade.

Mas não essa honestidade no sentido de não roubar e coisas do tipo. Refiro-me à honestidade comigo mesmo, com meus sentimentos e com as pessoas ao meu redor.

Não, não sou um livro aberto. Não se trata disso. Mas me sinto mal quando omito algo de alguém de quem gosto.

E não, não conto tudo o que existe dentro de mim porque sei que as pessoas ao meu redor não teriam como lidar com isso.

Isso porque o que existe aqui dentro faz parte desta jornada, que é particular e só minha.

Ao longo dessa caminhada, criei métodos, maneiras de lidar com o que muitos chamam de demônios.

Eu apelidei ou nomeei, seja lá como for, de método cebola.

Camadas de proteção.

Alguns chamam isso de criar casca.

Acredito que seja por isso que é tão difícil nos abrirmos cem por cento com as pessoas. E quem se abre acaba se machucando.

Assim, as camadas funcionam como amortecedores.

Com esse entendimento, adquiri a capacidade de ouvir as pessoas.

Sem julgamentos.

Sem estabelecer parâmetros.

Apenas ouvia e tentava compreender a perspectiva da pessoa, o que ela havia vivido e o que representava aquilo que estava falando.

E ficava me perguntando, enquanto ouvia: quantas cascas essa pessoa criou?

Passei a ser a pessoa confiável do grupo.

E mantinha em segredo o que me contavam, mesmo quando não me pediam isso.

Como isso se relaciona, segundo minha análise, sem nenhum embasamento teórico, com a parede branca?

No dia 11 de julho de 2026, veio a resposta, após um acontecimento, que me fez refletir.

Parei de falar sobre mim há mais de dez anos.

No meu egoísmo, eu pensava que as minhas coisas eram magníficas e de muita importância, e que todos tinham que saber, conhecer e se interessar por elas.

Então, passei a acreditar que ninguém se importava com o que eu falava e, consequentemente, não se importaria com o que eu colocasse em uma parede.

Já que me tornei um ouvinte, e não mais um falante.

Com o passar do tempo, a parede passou a representar o que eu sentia internamente: um vazio.

Uma falta de vontade de compartilhar.

Muitas coisas importantes que eu fazia já não tinha vontade de mostrar, de colocar em algum lugar onde outras pessoas pudessem ver.

E isso, não sei se justifica, mas pode ser um dos motivos da sensação de falta de propósito.

Não que em algum momente não tenha tentado falar, me abrir, inclusive sobre a parede branca e sobre o que ela despertou em mim, dada a minha história.

Mas eu era interrompido, e as pessoas começavam a falar sobre elas mesmas.

Acredito que assim a gente vai se fechando, e a parede vai ficando branca.

Como comentei no início, não se trata de um livro de autoajuda.

Logo, não há uma lição de moral, um aprendizado. Digo mais: não há um final feliz nem uma conclusão.

As pessoas são o que são.

Com suas limitações, seus medos, suas dores, suas alegrias, suas conquistas e suas derrotas.

A vida vai continuar seu ciclo, com começos e recomeços.

A parede pode até vir a ser preenchida, mas não é isso que vai mudar a vida real.

O fato é que estamos vivendo em um mundo onde as coisas que aparentemente têm valor são aquelas compartilhadas e exibidas nas redes sociais.

As pessoas nem se reúnem mais. Então, ter uma parede cheia ou vazia não vai fazer diferença.

Não se trata de um apego ao passado, até porque vocês viram meu relato sobre minha infância, adolescência e juventude.

Há coisas boas? Sim, há, lógico que há. Mas não gostaria de reviver aquela época.

Se você tem uma parede branca, não se sinta sozinho.

Se sente que não se encaixa no mundo de hoje, não se sinta sozinho.

Pena não haver um estudo sobre isso, mas é bem possível que muitas pessoas hoje estejam simplesmente lutando para se encaixar.

Então, para aqueles que estão nessa luta, só digo: continuem.

Um ciclo de cada vez.

Não se cobre tanto.

E não se compare com os outros.

Às vezes, a felicidade está em uma roda de amigos, com música, depois de um churrasco à beira da piscina.

Não estou sugerindo que você não tenha planos, sonhos, metas e objetivos.

Eu estou, neste momento, realizando um sonho: escrever um livro.

Algo que sempre achei que seria impossível.

Estou apenas sugerindo que não encare tudo o que acontece como se fosse algo definitivo e sem solução, nem subestime o poder do tempo.

Não subestime tudo o que já viveu.

E não se feche para o que ainda está por vir.

E que a vida apenas seja, que a vida apenas aconteça.